\chapter{A Worked Mathematical Example}
\label{chap:worked-example}


This chapter illustrates how to introduce notation, state a result, and connect
its proof to a numbered equation. The example concerns conjugacy classes in a
finite group. For research on weights and representations of finite groups,
see \citet{alperin1987weights} and \citet{knorr1989some}.


\section{Notation}
\label{sec:group-notation}
Let $G$ be a finite group. Write $k(G)$ for its number of conjugacy classes.
For $g \in G$, define its centralizer by
\[
  C_G(g) = \{h \in G : hg = gh\}.
\]
The conjugacy class of $g$ consists of the elements $hgh^{-1}$ for $h \in G$.
Its size is $|G| / |C_G(g)|$.


\section{Counting Conjugacy Classes}
\label{sec:conjugacy-classes}
\begin{prop}
  \label{prop:class-count}
  The number of conjugacy classes satisfies
  \begin{equation}
    \label{eq:class-count}
    k(G) = \frac{1}{|G|} \sum_{g \in G} |C_G(g)|.
  \end{equation}
\end{prop}
\begin{proof}
  Centralizers of elements in the same conjugacy class have the same size.
  A class represented by $g$ therefore contributes
  \[
    \frac{|G|}{|C_G(g)|} |C_G(g)| = |G|
  \]
  to the sum in Equation~\eqref{eq:class-count}. Summing over the $k(G)$
  conjugacy classes gives $k(G) |G|$. Dividing by $|G|$ proves the result.
\end{proof}
\begin{corollary}
  If $G$ is abelian, then $k(G) = |G|$.
\end{corollary}
\begin{proof}
  In an abelian group, $C_G(g) = G$ for every $g \in G$. Apply
  Proposition~\ref{prop:class-count}.
\end{proof}


For example, the cyclic group of order 3 has three conjugacy classes, each
containing one element. This example checks Equation~\eqref{eq:class-count}
without introducing additional notation.
